% ---------------------------------------------------------------------------
% 3 Deployment scenario and threat model
% Source: context/01_scenario_and_motivation.md  (this section is written FROM that note)
%         context/06_defensibility.md §3.1
% This is the spine. Sections 5, 6, and 7 are read against it.
% ---------------------------------------------------------------------------

\section{Deployment scenario and threat model}
\label{sec:threat-model}

\subsection{Deployment roles and privacy objective}
The deployment has two parties but one cryptographic confidentiality objective. The data owner
holds a human genomic sequence, ordinary CPU resources, and no GPU. It will not disclose the
sequence to the compute provider. The model owner operates the accelerator infrastructure and
serves the DNAGPT model rather than distributing it to clients. Keeping the weights server-side is
an operational preference of this deployment, not a model-confidentiality guarantee supplied by
the protocol.

The released DNAGPT artifact makes the experiment reproducible; it does not itself require this
service arrangement. A client may run the public release locally. We use it as an auditable
surrogate for a genomic model whose owner limits access to a service by policy or contract, and the
deployment conclusions apply under that premise.

The protected asset is therefore the genome. The data owner encrypts its embedded numeric vectors
and retains the secret key; the compute provider evaluates the served model on those ciphertexts.
The protocol is relevant only where these roles are fixed. If a client is allowed to receive and
run the model, local plaintext inference is the simpler architecture.

\subsection{Local inference as a comparison}
The sharpest objection to this protocol is best stated with our own measurements, at full strength,
before it is answered. The data owner already evaluates LayerNorm, causal softmax, and GELU, so it
is not spared the nonlinear work. In the complete encrypted inference reported in
Section~\ref{sec:results} it spends $956.187$~s of its own CPU time doing so, against a plaintext
forward pass of this released $0.1$-billion-parameter model at $103$ tokens that finishes in well
under a second on the same class of hardware. The data owner therefore performs orders of magnitude
more work inside the protocol than it would need to run the model itself, and in this study the
weights are public, so no model-secrecy argument compels the arrangement either. Every one of those
statements is correct, and none of them is contested below.

The answer is not that the arithmetic is wrong. It is that the local baseline is not a cheaper
route to the same answer but a different deployment, available only where the client may hold the
model in plaintext. Where model distribution is permitted, the client can avoid all cryptographic
work by running the ordinary forward pass, and it should. Client assistance cannot provide a
computational advantage over that option, and this paper does not claim one. Its purpose is to make
a service-only model available without sending the genomic input in plaintext.

That comparison changes the deployment premise. A served model is available to the data owner only
through the provider's interface; handing the client the weights creates a different service. Under
the premise studied here, the available choices are to disclose the genome, decline the inference,
or evaluate through a privacy-preserving protocol. Client-assisted CKKS addresses the last choice.
The reason to use it is access to a served model without genomic disclosure, not an attempt to beat
plaintext inference.

\begin{table}[t]\centering\small
\caption{Resource and trust allocation under local and served-model deployment.}
\label{tab:deployments}
\begin{tabularx}{\textwidth}{Y c c}
\toprule
 & Client-local inference & Client-assisted CKKS \\
\midrule
Client must hold the model & yes & no \\
Server observes the genome & not applicable & no \\
Client requires a GPU & deployment-dependent & no \\
Server evaluates model linear algebra & no & on ciphertexts \\
Client evaluates model nonlinearities & as part of local model & at declared boundaries \\
\bottomrule
\end{tabularx}
\end{table}

The protocol assigns encrypted dense linear algebra to the party with accelerator hardware while
the key holder evaluates exact nonlinearities on CPU. This allocation protects the genomic input
without requiring the data owner to operate a GPU.

This allocation becomes more server-heavy as model width grows. At fixed sequence structure, the
client's boundary work scales with the number of activation values, whereas the server's dense
transforms scale quadratically with hidden width. The client fraction therefore decreases as
$O(1/D)$ with hidden width $D$ under this fixed sequence structure. This is an asymptotic property
of the work allocation, not a measurement of a wider DNAGPT model.

\subsection{Adversary model}
We consider a semi-honest compute provider under $128$-bit classical CKKS parameters. It follows
the prescribed circuit but may inspect everything it legitimately receives: public and evaluation
keys, model parameters, ciphertexts, sequence length, message sizes, operation order, boundary
schedule, and timing. It does not alter ciphertexts to probe the client or deviate from the
declared boundary sequence.

Under this adversary, the compute provider receives no plaintext genomic vector, plaintext
activation, partial decryption, or secret-key material. At a boundary it supplies a ciphertext and
receives a fresh ciphertext under the data owner's public key. The guarantee concerns the
provider's view of query-derived values; public model structure and plaintext weights remain
available to the provider as inputs to ciphertext--plaintext operations.

\begin{invariant}[Key custody]
The data owner generates and retains the secret key. No protocol message contains the key or a
partial decryption, and the compute provider has no decryption operation.
\end{invariant}

\begin{invariant}[No plaintext at the server]
Every query-derived value held by the compute provider is a ciphertext under the data owner's key.
Client-evaluated nonlinearities return through fresh encryption rather than as decoded values.
\end{invariant}

\begin{invariant}[Non-adaptive boundaries]
The circuit and sequence length determine every boundary before execution; decrypted content does
not select the next operation. The implementation asserts the complete schedule and boundary counts
for each accepted run and aborts on deviation. Transcript equality across multiple private inputs
of equal length has not been measured as a separate leakage experiment.
\end{invariant}

\subsection{What is not claimed}
Model confidentiality is not claimed. The data owner observes full intermediate activations at
every nonlinearity. Those observations decouple the network into shallow segments with known
nonlinearities, turning global model inversion into inexpensive per-segment regression. Keeping
weights on the server may deter casual copying, but it does not prevent extraction by a determined
client. This concession does not change the genomic-data guarantee, which protects the data owner
from the compute provider.

Malicious servers, authenticated transport, production key custody, traffic analysis, timing and
physical side channels, and compromised clients are outside the evaluated boundary. Encrypted
scope begins after token-index lookup, so private lookup is also unresolved. The embedded vectors that cross the boundary are not a weaker secret than the sequence: embeddings from DNA foundation models have been inverted to recover the underlying sequence with high fidelity~\cite{dnaembeddinginversion2026}, which is why they are encrypted here rather than sent in clear. The data owner sees
its own intermediate activations and final output; privacy from the key holder is not an objective
of this protocol.
